\subsection{Entrevista}

En esta sección se detallan las entrevistas realizadas a los diferentes actores del consorcio para comprender sus necesidades y puntos de vista respecto al funcionamiento del sistema de administración.

\vspace{0.3cm}
\textbf{Entrevista 1: Administrador del Consorcio}

Se utilizó una estructura de embudo para la entrevista hacia el administrador del consorcio.  

\textbf{Preguntas abiertas:}

1. ¿Cuáles son los objetivos críticos de su departamento al administrar este edificio?

\textbf{Administrador:} ``El objetivo principal es mantener las cuentas al día para poder pagar los sueldos y los servicios sin sobresaltos. Además, necesito asegurar que el edificio esté en buenas condiciones de mantenimiento y, sobre todo, transparentar los gastos para que los vecinos no se quejen todo el tiempo. Básicamente, que la plata alcance y tener paz en el consorcio.'' 

2. Describa el proceso actual para calcular y distribuir las expensas mensuales.

\textbf{Administrador:} ``Hoy es un dolor de cabeza tremendo. A fin de mes, junto todos los tickets de gastos (luz, agua, arreglos, sueldo del encargado), los sumo en una planilla de Excel, divido por los porcentajes de cada departamento según los metros cuadrados, y armo un PDF. Después, tengo que mandarlos uno por uno por mail o imprimir los papelitos para pasarlos por debajo de la puerta. Me lleva días.''   

\textbf{Pregunta de Sondeo:} 

3. ¿Me lo puede explicar con más detalle para los casos en que hay deudores morosos? 

\textbf{Administrador:} ``Uf, ese es el mayor problema. En el Excel tengo que ir marcando quién pagó y quién no. Si alguien se atrasa, al mes siguiente tengo que acordarme de sumarle el interés punitorio del 5\% a mano sobre su deuda anterior, más la expensa nueva. A veces se me pasa, o el vecino me dice 'yo te transferí el día 15' y tengo que ir a buscar en el resumen del banco línea por línea a ver si es verdad. Es muy poco práctico.''   

\textbf{Preguntas cerradas:}

4. Mencione sus dos prioridades principales para mejorar la infraestructura de tecnología del edificio.

\textbf{Administrador:} ``Primero, automatizar el cálculo y el envío de las expensas. Segundo, tener un registro automático de los pagos para no tener que cruzar las transferencias bancarias a mano todos los días.''   

¿Cuántos subordinados tiene a cargo para el mantenimiento y seguridad que usarán el sistema?  

\textbf{Administrador:} ``Directamente relacionados con el uso del sistema, tengo a los dos guardias de seguridad que rotan en la entrada y al encargado del edificio. Tres en total.''   

\textbf{Preguntas bipolares:} 

5. ¿Desea recibir una impresión de su estado de cuenta cada mes?

\textbf{Administrador:} ``No, prefiero que quede todo 100\% digital e histórico en el sistema para poder consultarlo desde la computadora cuando quiera, sin acumular más papel.'' 

6. ¿Está de acuerdo o en desacuerdo con que el actual método manual de reservas de los espacios comunes carece de seguridad?

\textbf{Administrador:} ``Totalmente de acuerdo. Hoy con el cuaderno de papel que tiene el guardia, cualquiera tacha un nombre o se anotan dos vecinos en el mismo horario para usar los espacios comunes, y después el problema y las discusiones me explotan a mí.'' 

\vspace{0.5cm}
\textbf{Entrevista 2: Propietario/Cliente}

Para esta entrevista utilizamos la estructura de pirámide, ya que es la más acorde para el cliente.  

\textbf{Preguntas bipolares:}

1. ¿Actualmente recibe el resumen de sus expensas impreso en papel por debajo de su puerta? 

\textbf{Cliente:} ``Sí, tal cual. Me lo tiran por debajo de la puerta todos los meses.''  

\textbf{Preguntas cerradas:}

2. ¿Cuántas veces al año suele reservar los espacios comunes, como el quincho o los espacios comunes? 

\textbf{Cliente:} ``Calculale unas tres o cuatro veces al año como mucho, más que nada para algún cumpleaños o una cena familiar grande.''  

\textbf{Pregunta de Sondeo:}

3. Mencionó que el método para reservar es anotarse en un cuaderno en la guardia, ¿me puede dar un ejemplo de qué ocurre si ese cuaderno se pierde o hay doble reserva? 

\textbf{Cliente:} ``Uy, es un desastre. Una vez  organisé el cumple de mi nene, compré la comida, invité a todos, y cuando bajé al quincho con las cosas había otra familia instalada. Resulta que el guardia del turno mañana me anotó en una hoja, pero el del turno noche no lo vio y le reservó a otro vecino en el mismo horario. Nos terminamos peleando entre nosotros por culpa de la desorganización.''  

\textbf{Preguntas abiertas:}

4. ¿Cuáles son los métodos de pago electrónico que preferiría tener integrados en la nueva plataforma? 

\textbf{Cliente:} ``Para mí sería clave que tenga MercadoPago o algún QR para pagar en el momento desde el celular. Y si es transferencia bancaria, que el sistema te lo tome automático, porque hoy en día tenés que hacer la transferencia y después andar mandándole un mail al administrador con la captura de pantalla para que te lo anote. Es una pérdida de tiempo.''

5. En general, ¿qué opina de la transparencia actual en el detalle de los gastos y cómo cree que el sistema podría mejorar esa comunicación? 

\textbf{Cliente:} ``La verdad, hoy la transparencia es nula. Te llega esa hoja por debajo de la puerta con un número final gigante y un resumen que no se entiende nada. A veces dice 'Arreglo de caños: \$200.000' y vos no tenés idea de qué caño se rompió ni si pidieron presupuesto. Yo creo que si el sistema tuviera una sección tipo 'Muro' o de 'Documentos' donde suban las facturas de lo que se arregla, o un gráfico para ver en qué se nos va la plata mes a mes, nos daría muchísima más tranquilidad a todos.''
